\scanpage{257}
\[
\varinjlim\,\unsure{\varphi|\rho}_N\tau_N P\longrightarrow\varphi_\infty\tau_\infty P
\]
\deleted{Hom}\quad bijectif ?
\[
\Downarrow
\]
$\forall M$ de \unsure{prés|prés.} finie (il suffit $M=S[i]$)
\diagram{\begin{tikzcd}[ampersand replacement=\&,column sep=small,execute at end picture={\draw ([xshift=-2pt]\tikzcdmatrixname-2-2.west |- \tikzcdmatrixname-2-2.north) -- ([xshift=-2pt]\tikzcdmatrixname-2-2.west |- \tikzcdmatrixname-4-1.south);}]
\operatorname{Hom}_{\mathcal M}(M,\underset{\unsure{\rightarrow|-}}{\lim}\varphi_N\tau_N P)\arrow[r,dash]\arrow[d,phantom,"\rotatebox{90}{$\simeq$}"] \& \operatorname{Hom}_{\mathcal M}(M,\varphi_\infty\tau_\infty P)\quad\text{bij}\arrow[d,phantom,"\rotatebox{90}{$\simeq$}"]\\
\varinjlim\operatorname{Hom}_{\mathcal M}(M,\deleted{\unsure{\tau}}\varphi_N\tau_N P)\arrow[d,phantom,"\rotatebox{90}{$\simeq$}"] \&\operatorname{Hom}_{\mathcal M_\infty}(\tau_\infty M,\tau_\infty P)\arrow[d,phantom,"\rotatebox{90}{$\simeq$}"]\\
\varinjlim\operatorname{Hom}_{\mathcal M_N}(\tau_N M,\tau_N P)\arrow[d,phantom,"\rotatebox{90}{$\simeq$}"] \& \underset{\substack{\longrightarrow\\ M'\subset M\\M/M'\ \text{négl.}}}{\lim}\operatorname{Hom}(M',P)\\
\underset{\unsure{-|\rightarrow}}{\lim}\operatorname{Hom}_{\mathcal M}(\tau_N M,P)\arrow[ur]\&
\end{tikzcd}}

dire que (pour $M$ fixé) c'est un isom pour tt $P$, signifie que les sous pro-objets de $M$ \unsure{formés}, soit des $\tau_N M$, soit des $M'\subset M$ tels que $M/M'$ soit négligeable, \deleted{\unsure{sont|soit}} isom par l'hom can., ou encore que, \unsure{ces ens|ces deux} de sous-objets sont cofinaux, i.e que tout $M'$ contient un $\tau_N M$ \inserted{i.e $M/M'$ est $\infty$-négligeable.}

Mais c'est le cas si $M$ de \deleted{\unsure{prés.}}\inserted{t.f.}

\deleted{\uline{Hyp} Tout idéal gradué $J$ de $S$, tel que $S/J$ soit négligeable}

\deleted{\unsure{finité}.} Donc
\[
\varinjlim\,\varphi_N\tau_N P\xrightarrow{\sim}\varphi_\infty\tau_\infty P
\]
% archive sheet number: 256
